\section{Compactness and quadratic envelopes}\label{sec:limiting-tools}

We begin the proof of Theorem~\ref{thm:liouville}. Write
$\Aclass_{m,C_0}$ for the class of continuous functions on
$\R^m\times\R$ that solve \eqref{eq:limit-pde} and have the fits
\eqref{eq:limit-fit} at every center and radius. We may enlarge $C_0$ to
$\max\{C_0,1\}$. Constants denoted by $C$ in estimates for this class
may depend on $m$ and $C_0$.

\subsection{Changes of frame and compactness}

\begin{lemma}[Changes of frame]\label{lem:frame-invariance}
Let $v\in\Aclass_{m,C_0}$, $z_0=(x_0,t_0)$, $R>0$, and
$(p_0,a_0)\in\R^m\times\R$. Then
\begin{equation}\label{eq:frame-change}
  \bar v(y,s)=
  \frac{v(x_0+R(y+p_0s),t_0+Rs)-v(x_0,t_0)}{R}
       -p_0\cdot y-\left(a_0+\frac{|p_0|^2}{2}\right)s
\end{equation}
belongs to $\Aclass_{m,C_0}$. A fit $(p,a)$ for $v$ at the
corresponding center and radius $R\rho$ becomes the fit
$(p-p_0,a-a_0)$ for $\bar v$ at radius $\rho$.
\end{lemma}

\begin{proof}
Set
\[
  B=\begin{pmatrix}I_m&p_0\\0&1\end{pmatrix}.
\]
For smooth functions, at corresponding points,
$\nabla_y\bar v=\nabla_xv-p_0$ and
$D^2\bar v=R B^{\mathsf T}D^2vB$. Since
$B(\nabla_y\bar v,1)=(\nabla_xv,1)$, this gives
$F[\bar v]=R F[v]$. Pulling back an upper or lower smooth test gives
the same implication in the viscosity sense. For the fits, an
increment $(h,\ell)$ in the new variables corresponds to
$R(h+p_0\ell,\ell)$. The old affine increment, divided by $R$ and
with the affine part of \eqref{eq:frame-change} subtracted, is
\[
  (p-p_0)\cdot h+
  \left(a-a_0-\frac{|p-p_0|^2}{2}\right)\ell.
\]
The old box conditions become
$|\ell|\le\rho$ and $|h-(p-p_0)\ell|\le\rho$, and the error is
$C_0\rho$.
\end{proof}

We say that $v$ is \emph{centered at $z_c$} if $v(z_c)=0$ and
$(0,0)$ is a fit there of radius $1$. Choosing $(p_0,a_0)$ to be a
fit at $(z_0,R)$ in \eqref{eq:frame-change} produces a function
centered at the origin. We will also center at $(0,1)$, using a
translation of the time coordinate. Put
\[
  L(K)=1+\log K\qquad(K\ge2),\qquad
  \mathcal B_K=\{(x,t):|x|\le K,\ |t|\le K\}.
\]

Comparisons of fits across scales give the logarithmic parameter bounds
below. They yield compactness after centering and uniform subquadratic
spatial growth, which will confine the quadratic-envelope optimizers.
The estimates use only the fits; the equation enters when the limit is
identified.

\begin{lemma}[Bounds and compactness from fits]\label{lem:fit-compactness}
Suppose that $v\in\Aclass_{m,C_0}$ is centered at either $(0,0)$ or
$(0,1)$. For $K\ge2$ one has
\begin{equation}\label{eq:centered-growth}
  |v(x,t)|\le C K L(K)^2\qquad ((x,t)\in\mathcal B_K).
\end{equation}
Every fit $(p,a)$ at a center in $\mathcal B_K$, of any radius $r>0$,
satisfies
\begin{equation}\label{eq:fit-parameter-bound}
  |p|+|a|\le C\bigl(L(K)+|\log r|\bigr).
\end{equation}
In particular, for a center $(x,t)\in\mathcal B_K$,
\begin{align}
  |v(x+h,t)-v(x,t)|
    &\le Cr\bigl(L(K)+|\log r|\bigr)
       &&(|h|\le r),\label{eq:spatial-fit-modulus}\\
  |v(x,t+\tau)-v(x,t)|
    &\le C|\tau|\bigl(L(K)+|\log|\tau||\bigr)^2
       &&(0<|\tau|\le\tfrac12).\label{eq:time-fit-modulus}
\end{align}
For either fixed choice of center, every sequence of centered
members has a locally uniformly convergent subsequence, and its
limit belongs to $\Aclass_{m,C_0}$ and is centered at that same
point. Moreover, on every compact time interval these members have
uniformly subquadratic growth in $x$:
\begin{equation}\label{eq:uniform-subquadratic-growth}
  \lim_{M\to\infty}\ 
  \sup_{\substack{v\text{ centered}\\ |x|\ge M,\ t\in J}}
       \frac{|v(x,t)|}{|x|^2}=0
  \qquad(J\Subset\R).
\end{equation}
\end{lemma}

\begin{proof}
We prove the estimates for centering at the origin. Translation by
one unit in time gives the other case after enlarging the boxes by a
fixed factor.

First compare two fits $(p,a),(p',a')$ at the same center, with
radii $R$ and $R'\in[R,2R]$. On the horizontal slice $\ell=0$,
both estimates hold for $|h|\le R$, so
$|p-p'|\le 3C_0\le4C_0$. Write
\[
  \ell_{p,a}(h,\ell)=p\cdot h+
                \left(a-\frac{|p|^2}{2}\right)\ell.
\]
With $c=(2+8C_0)^{-1}$, the point $(p cR,cR)$ lies in both sheared
boxes. At this point
\[
  (\ell_{p,a}-\ell_{p',a'})(p cR,cR)
    =\left(a-a'+\frac{|p-p'|^2}{2}\right)cR.
\]
The two fitting errors sum to at most $3C_0R$. It follows that
$|a-a'|\le 3C_0/c+8C_0^2$. Thus the differences of both parameters
are bounded by one constant whenever the radii differ by at most a
factor of two. Comparing along successive dyadic radii with the
centered fit of radius $1$ gives, for every fit at the origin of
radius $r$,
\begin{equation}\label{eq:origin-fit-bound}
  |p|+|a|\le C(1+|\log r|).
\end{equation}

Fix $K\ge2$. Choose once and for all a sufficiently large constant
$A=A(m,C_0)$ and set
$R=A K L(K)$. If $(p_R,a_R)$ is any fit of radius $R$ at the
origin, then $\mathcal B_K$ lies in its sheared box with
\[
  |t|\le R/4,\qquad |x-p_Rt|\le R/4.
\]
Indeed, \eqref{eq:origin-fit-bound} bounds the latter offset by
$K(1+C(1+\log R))$, while
$\log R\le\log A+2L(K)$. Since $L(K)\ge1$, the asserted two
inequalities hold for all $K\ge2$ once $A$ is chosen larger than a
fixed multiple of $1+\log A$. Such a choice is possible because
$\log A/A\to0$. The fit at radius $R$ now gives
\[
  |v(x,t)|
  \le C_0R+|p_R|K+
      \left(|a_R|+\frac{|p_R|^2}{2}\right)K
  \le C K L(K)^2,
\]
which is \eqref{eq:centered-growth}.

Let $z=(x,t)\in\mathcal B_K$, and compare a fit $(p_z,a_z)$ of
radius $R$ at $z$ with the enclosing fit $(p_R,a_R)$. For
horizontal increments of length at most $R/2$ from $z$, the
enclosing estimate can be subtracted at the two endpoints, with
error at most $2C_0R$. Comparison with the fit at $z$ therefore
gives $|p_z-p_R|\le6C_0$. Set
$c'=(4(1+6C_0))^{-1}$. The increments
$(p_z c'R,c'R)$ from $z$ remain in both boxes: in the enclosing
box the transverse offset is at most
$R/4+6C_0c'R\le R/2$, and the time offset is at most $R/2$.
The same affine-increment identity used above gives a bound on
$|a_z-a_R|$ depending only on $C_0$. Hence
$|p_z|+|a_z|\le C L(K)$. Dyadic comparisons at $z$, from radius
$R$ to an arbitrary radius $r$, give
\[
  |p|+|a|
  \le C L(K)+C(1+|\log(r/R)|)
  \le C\bigl(L(K)+|\log r|\bigr),
\]
as claimed. These comparisons apply to every choice of fit.

Using a fit of radius $r$ on the horizontal slice proves
\eqref{eq:spatial-fit-modulus}. For a time increment $\tau$ with
$0<|\tau|\le1/2$, put
$B=L(K)+|\log|\tau||$ and choose the fit radius
$r=A'|\tau|B$. Here $A'$ is a constant to be fixed, independent of
$K$ and $\tau$. Since $B\ge1$,
\[
  |\log r|\le|\log|\tau||+\log A'+\log B
       \le2B+\log A'.
\]
The fit-parameter bound consequently gives
$|p|\le C(3+\log A')B$. Choose $A'\ge1$ sufficiently large that
$A'\ge C(3+\log A')$. This choice is possible and is made before
$K$ and $\tau$ vary. Then $|\tau|\le r$ and $|p\tau|\le r$, so
the vertical increment lies in the fitted box. In that estimate,
the affine time coefficient and error are bounded by
\[
  \left(|a|+\frac{|p|^2}{2}\right)|\tau|+C_0r
      \le C|\tau|B^2.
\]
This proves \eqref{eq:time-fit-modulus}.

The two moduli, applied successively to a spatial and a time
increment in a slightly larger box, give equicontinuity on compact
sets. The growth bound gives local uniform boundedness.
Arzel\`a--Ascoli and a diagonal choice of subsequence therefore
give local uniform convergence. To pass an upper viscosity test
to the limit, add a positive fourth-order term at its contact,
preserving its first and second derivatives and making the
contact strict on a small closed ball. Maxima of the sequence
minus this modified test on that ball converge to the contact,
so they are interior for large indices. The viscosity
inequalities there pass to the limit by continuity of $F$.
Subtracting the fourth-order term and using minima treats a
lower test. To pass the fits, fix one
center and one radius $r$. By \eqref{eq:fit-parameter-bound}, pass
to a further subsequence of the corresponding fits with
$(p_j,a_j)\to(p,a)$. For a point with
$|\ell|\le r$ and $|h-p\ell|\le r$, use in the $j$th fitted box
the increment
$h_j=p_j\ell+(h-p\ell)$. It converges to $h$ and belongs to that
closed box. Local uniform convergence gives the limiting
inequality. This subsequence may depend on the fixed center and
radius, which is sufficient for existence of a fit at each such
pair. The centered fit passes directly to the limit.
Finally, for $t$ in a fixed compact interval,
\eqref{eq:centered-growth} gives a bound of the form
$C_J(1+|x|)(1+\log(2+|x|))^2$, uniformly in the centered class.
Dividing by $|x|^2$ proves
\eqref{eq:uniform-subquadratic-growth}.
\end{proof}

If a family instead has one common prescribed value and one common
fit at a fixed center and radius, one fixed change of frame
\eqref{eq:frame-change} reduces it to the centered case. The
corresponding local bounds and compactness hold with constants
also depending on those prescribed data.

\subsection{Quadratic envelopes and their slopes}

Let $I\subseteq\R$ be an open interval. We say that a continuous
function $f$ on $\R^m\times I$ has \emph{subquadratic spatial
growth on compact time intervals} if, for every $J\Subset I$,
\[
  \lim_{M\to\infty}
  \sup_{\substack{|x|\ge M\\t\in J}}\frac{|f(x,t)|}{|x|^2}=0.
\]
For $\ell>0$, with the indicated time argument in $I$, define
\begin{equation}\label{eq:envelope-definitions}
  \begin{split}
  \Penv_\ell f(x,t)
    &=\sup_y\left\{f(y,t+\ell)-\frac{|y-x|^2}{2\ell}\right\},\\
  \Menv_\ell f(x,t)
    &=\inf_y\left\{f(y,t-\ell)+\frac{|y-x|^2}{2\ell}\right\}.
  \end{split}
\end{equation}
Set $\Penv_0f=\Menv_0f=f$. The following quadratic slice
comparisons are useful for the limiting equation; compare
\cite[Section~3]{EvansSavin2008}. Convexity of Hamilton--Jacobi envelopes also
appears in the comparison theory of Armstrong, Crandall, Julin, and Smart
\cite[Proposition~2.3]{ArmstrongCrandallJulinSmart2011}. Their Hamiltonian
hypotheses do not cover the limiting equation here; we prove the required
properties directly.
For a fixed center, convexity in the length will give monotone envelope
slopes. The moving-center form records the cost of shifting the center;
its midpoint inequality will control velocity differences in
Section~\ref{sec:discrete-propagation}.

\begin{lemma}[Quadratic envelopes]\label{lem:quadratic-envelopes}
Suppose that $f$ has the preceding growth property. The extrema
in \eqref{eq:envelope-definitions} are attained. The envelopes
are continuous jointly in their center and length, including
length zero wherever the time arguments remain in $I$. Whenever
the expressions are defined,
\begin{equation}\label{eq:envelope-composition}
  \Penv_\ell(\Penv_hf)=\Penv_{\ell+h}f,\qquad
  \Menv_\ell(\Menv_hf)=\Menv_{\ell+h}f.
\end{equation}
If $f$ is a viscosity subsolution of \eqref{eq:limit-pde}, then,
for every affine path $x(\ell)=x_0+d\ell$ and every fixed $t$,
\begin{equation}\label{eq:path-envelope-convexity}
  \ell\longmapsto
  \Penv_\ell f(x(\ell),t)+|d|^2\ell\log\ell
  \quad\text{is convex}
\end{equation}
on its allowed range $\ell\ge0$, with $\ell\log\ell=0$ at zero.
In particular $\ell\mapsto\Penv_\ell f(x,t)$ is convex.
If $f$ is a viscosity supersolution, then
$\ell\mapsto\Menv_\ell f(x,t)$ is concave.
If $f_j\in C(\R^m\times I)$, $f_j\to f$ locally uniformly, and
\[
  \lim_{M\to\infty}\sup_j\sup_{\substack{|y|\ge M\\t\in J}}
      \frac{|f_j(y,t)|}{|y|^2}=0\qquad(J\Subset I),
\]
then
$\Penv_\ell f_j\to\Penv_\ell f$ and $\Menv_\ell f_j\to\Menv_\ell f$
locally uniformly in the allowed centers and positive lengths.
\end{lemma}

\begin{proof}
Fix a compact range of centers and of allowed positive lengths.
Uniform subquadratic growth and the quadratic cost force all
maximizers and minimizers into a common bounded spatial set.
They exist there, and the maximum and minimum of a continuous
function over this fixed compact set depend continuously on the
parameters. The same bounded-set conclusion holds when the
length varies in $(0,\ell_0]$. On that set the values of $f$ are
bounded. Comparison with the competitor $y=x$ then gives
$|y-x|^2\le C\ell$ for an optimizer as $\ell\downarrow0$.
Uniform continuity of $f$ on the compact set gives convergence
to $f(x,t)$ locally uniformly, as well as $y\to x$ locally
uniformly. This proves continuity at zero. The composition
identities follow by optimizing first over the intermediate
point and using
\[
  \inf_y\left\{\frac{|y-x|^2}{2\ell}
           +\frac{|z-y|^2}{2h}\right\}
       =\frac{|z-x|^2}{2(\ell+h)}.
\]

For the convexity assertion, let
$g(\ell)=\Penv_\ell f(x(\ell),t)$, and let a smooth function
$\psi$ touch $g$ from above at $\ell_0>0$. Choose an optimizer
$y_0$ at that length. After a constant adjustment,
\[
  \Psi(y,\ell)=\psi(\ell)+\frac{|y-x(\ell)|^2}{2\ell}
\]
touches $(y,\ell)\mapsto f(y,t+\ell)$ from above at
$(y_0,\ell_0)$. Put $q=(y_0-x(\ell_0))/\ell_0$, which is the
horizontal gradient of $\Psi$ there. Along the direction
$(q,1)$, the Hessian of the quadratic term has value
$|d|^2/\ell_0$. For example, writing $r=y-x(\ell)$, its value
along a direction with $r$-velocity $q-d$ is
$\ell_0^{-1}|q-d-r/\ell_0|^2=|d|^2/\ell_0$.
The viscosity inequality gives
$\psi''(\ell_0)+|d|^2/\ell_0\ge0$. Thus every upper test for
$g(\ell)+|d|^2\ell\log\ell$ has nonnegative second derivative
at positive lengths. This is the one-dimensional viscosity
criterion for convexity. To see it directly, a failure of
convexity on $[a,b]$ lets one place a strictly concave quadratic
above the chord but below the function somewhere; the positive
maximum of their difference yields an upper test with negative
second derivative. Continuity at zero extends the convexity to
that endpoint.

Finally, the map $\widetilde f(x,t)=-f(x,-t)$ exchanges
subsolutions and supersolutions of \eqref{eq:limit-pde}, as
pullback reverses the sign of $F$. It also gives
$\Penv_\ell\widetilde f(x,t)=-\Menv_\ell f(x,-t)$.
The assertion for $\Menv$ follows.

For the last convergence assertion, fix a compact range of allowed centers
and positive lengths. The common subquadratic bound and the local uniform
bounds confine every optimizer for $f_j$ and for $f$ to one compact spatial
set, by the first argument of the proof. On that set the functions inside
the supremum and infimum converge uniformly. Their extrema therefore
converge uniformly on the chosen parameter range.
\end{proof}

For a positive allowed length define the two slopes
\begin{equation}\label{eq:envelope-slopes}
  \Senv_\ell^+f(z)=\frac{\Penv_\ell f(z)-f(z)}{\ell},
  \qquad
  \Senv_\ell^-f(z)=\frac{f(z)-\Menv_\ell f(z)}{\ell}.
\end{equation}
For a smooth function, the quantity
$H_{\mathrm{sm}}=f_t+|\nabla_xf|^2/2$ satisfies
\[
  F[f]=(\partial_t+\nabla_xf\cdot\nabla_x)H_{\mathrm{sm}}.
\]
For the exact affine model
$f(x,t)=p\cdot x+(a-|p|^2/2)t+c$ in \eqref{eq:limit-fit}, one has
$H_{\mathrm{sm}}=a$, and direct optimization gives
\[
  \Penv_\ell f=f+a\ell,\qquad \Menv_\ell f=f-a\ell.
\]
The optimizing displacements are respectively $\ell p$ and $-\ell p$.
Thus both envelope slopes equal $a$, and both optimizer velocities,
directed toward increasing time, equal $p$. These identifications concern
the exact affine model; for a general fitted solution, $p,a$ remain the
parameters of the chosen fit. For a general smooth function, a first-order
expansion of the two quadratic optimizations formally gives
$\Senv_\ell^\pm f=H_{\mathrm{sm}}+o(1)$ as $\ell\downarrow0$.
For a smooth solution the identity for $F[f]$ transports this quantity along
the direction $(\nabla_xf,1)$; its value may depend on the integral curve.
The next lemma supplies the slope statements used for continuous solutions.

Evans and Savin \cite[Theorem~3.1]{EvansSavin2008} proved the corresponding
slope monotonicity and equality of the finite limits at a fixed center under
their local boundedness and Lipschitz assumptions. The following formulation
uses envelopes on the full spatial domain under subquadratic growth,
separates the subsolution and supersolution conclusions, and allows the common
limit to be $-\infty$.

\begin{lemma}[Envelope slopes]\label{lem:envelope-slopes}
For a subsolution $f$ as in Lemma~\ref{lem:quadratic-envelopes},
$\Senv_\ell^+f(z)$ is nondecreasing in $\ell$ and
\begin{equation}\label{eq:subsolution-slope-order}
  \limsup_{\ell\downarrow0}\Senv_\ell^-f(z)
       \le\lim_{\ell\downarrow0}\Senv_\ell^+f(z).
\end{equation}
The right side may equal $-\infty$. For a supersolution,
$\Senv_\ell^-f(z)$ is nondecreasing and
\[
  \limsup_{\ell\downarrow0}\Senv_\ell^+f(z)
       \le\lim_{\ell\downarrow0}\Senv_\ell^-f(z).
\]
Consequently, for a
solution both slopes are nondecreasing and
\begin{equation}\label{eq:hamiltonian}
  \lim_{\ell\downarrow0}\Senv_\ell^+f(z)
  =\lim_{\ell\downarrow0}\Senv_\ell^-f(z)
  =:H_f(z)\in[-\infty,\infty).
\end{equation}
The function $H_f$ is upper semicontinuous.
\end{lemma}

\begin{proof}
Convexity of $\Penv_\ell f(z)$, with value $f(z)$ at zero,
gives the first monotonicity. Let
$\widetilde z_\ell=(y_\ell,t-\ell)$ be a backward optimizing
point for $\Menv_\ell f(z)$. Using $z$ as a forward competitor
from $\widetilde z_\ell$ gives
\[
  \Senv_\ell^-f(z)\le\Senv_\ell^+f(\widetilde z_\ell).
\]
The optimizer localization at length zero in
Lemma~\ref{lem:quadratic-envelopes} shows
$\widetilde z_\ell\to z$. Fix a small $h>0$ whose forward time
argument is allowed near $z$. For all sufficiently small
$\ell<h$, monotonicity bounds the right side by
$\Senv_h^+f(\widetilde z_\ell)$. Continuity for fixed $h$ gives
an upper limit at most $\Senv_h^+f(z)$. Now let $h\downarrow0$
to obtain \eqref{eq:subsolution-slope-order}. This order of
limits also covers the value $-\infty$.

Apply this conclusion to $\widetilde f(x,t)=-f(x,-t)$ to obtain
the supersolution statements. For a solution both monotone
limits exist and the two inequalities identify them. Each is
bounded above by the finite slope at any fixed positive allowed
length, proving the range in \eqref{eq:hamiltonian}. Locally in
$z$, choose a common small range of lengths. The common limit is
the infimum of the continuous functions $\Senv_\ell^+f(z)$ over
that range, so it is upper semicontinuous.
\end{proof}

\begin{lemma}[Bounds for fitted solutions]\label{lem:fitted-envelope-bounds}
For a member $v\in\Aclass_{m,C_0}$ centered at $(0,0)$ or $(0,1)$, $K\ge2$,
$z\in\mathcal B_K$, and $0<\ell\le K$, one has
\begin{equation}\label{eq:envelope-slope-bound}
  \Senv_\ell^+v(z)\le C L(K)^4,\qquad
  \Senv_\ell^-v(z)\le C L(K)^4.
\end{equation}
Every optimizing displacement for $\Penv_Kv(z)$ or
$\Menv_Kv(z)$ has length at most $C K L(K)^2$.
\end{lemma}

\begin{proof}
By Lemma~\ref{lem:envelope-slopes}, it suffices to bound the two
slopes at length $K$. Suppose an optimizing displacement has
length $BK$ with $B\ge2$. Comparing its value with the
competitor having no displacement, and using
\eqref{eq:centered-growth} in a box of size $(B+2)K$, gives
\[
  \frac{B^2K}{2}
     \le C BK\bigl(L(K)+\log B\bigr)^2.
\]
This estimate holds for the supremum and for the infimum.
Consequently
$B\le C(L(K)+\log B)^2\le2C L(K)^2+2C(\log B)^2$.
Choose a fixed $B_0$ so large that
$2C(\log B)^2\le B/2$ for $B\ge B_0$. The last inequality
then yields $B\le4C L(K)^2$ for $B\ge B_0$, and the same type
of bound holds when $B<B_0$. This proves the displacement
claim. The optimizer and its time argument now lie in a box of
size $C K L(K)^2$. Applying \eqref{eq:centered-growth} there
and observing that $L(C K L(K)^2)\le C L(K)$ bounds the
absolute values of $v$ there by $C K L(K)^4$. Dropping the
nonpositive cost in the expressions for the two slopes gives
\eqref{eq:envelope-slope-bound}.
\end{proof}

The criterion for the final limit is proved in
Lemma~\ref{lem:slope-rigidity}: if $v\in\Aclass_{m,C_0}$ and there are
$(b,t_0)\in\R^m\times\R$ and $B\in\R$ such that
\[
  \Senv_\ell^+v(b,t_0)=\Senv_\ell^-v(b,t_0)=B
       \qquad\text{for every }\ell>0,
\]
then $v$ is affine. The remaining argument will construct such a limit while
retaining an obstruction to affinity.
